% TOP_PROBLEM: {"id": 30002825, "title": "Chromatic Fiber Structure at Bad Primes", "category_id": 7, "category": {"id": 7, "name": "topology", "display_name": "Topology", "description": "Properties preserved under continuous deformations.", "slug": "topology", "order_index": 7, "created_at": "2026-07-31T15:26:25.671Z"}, "status": "open"}
% FIRST_POSTED: 2026-10-06
% ATTEMPT_STATUS: unresolved
% OPENAI_CLAIM: {"schema_version":1,"source_repo":"https://github.com/openai/math","source_commit":"adc7f1241b42e322a6451854ab7e4b4c146bf78a","release_date":"2026-10-06","match":"partial","resolution":"partial","independently_reviewed":false,"second_pass":"n/a","families":[{"family":"318","title":"Chromatic splitting: filtrations and counterexamples","subject":"Topology","manuscripts":[{"title":"Counterexamples to weak chromatic splitting: sphere kernels and descent exponents","pdf_path":"preprints/Counterexamples-to-weak-chromatic-splitting-sphere-kernels-and-descent-exponents-September-27-2026/paper.pdf","readme_path":"preprints/Counterexamples-to-weak-chromatic-splitting-sphere-kernels-and-descent-exponents-September-27-2026/README.md","date":"2026-09-27","theorem_ref":"preprints/Counterexamples-to-weak-chromatic-splitting-sphere-kernels-and-descent-exponents-September-27-2026/build/sections/01-introduction.tex#thm:height-p-kernel"}],"lean":null,"reasoning_trace":null}]}
\documentclass[11pt]{article}
\usepackage[margin=25mm]{geometry}
\usepackage{fontspec}
\setmainfont{Latin Modern Roman}
\usepackage{amsmath,amssymb,mathtools}
\usepackage{unicode-math}
\setmathfont{Latin Modern Math}
\usepackage{enumitem}
\usepackage{xurl,hyperref,bookmark}
\hypersetup{colorlinks=true,urlcolor=blue}
\setcounter{secnumdepth}{0}
\setlength{\parindent}{0pt}
\setlength{\parskip}{5pt}
\setlength{\emergencystretch}{3em}
\begin{document}
\section*{Chromatic Fiber Structure at Bad Primes}

\subsection{OpenAI's claim}
The canonical weak-splitting unit is claimed to have no retraction at height p for p at least five and height p+1 for p at least seven, witnessed by \(\beta_1^{(p-1)^2}\). The cited manuscript argues for that result; its relation to this problem is specified below.

\subsection{Outline of the argument}
Manuscript 1 (family 318, 2026-09-27): The named beta power survives under the height-(p-1) unit by an ordinary homotopy-fixed-point detection calculation. At each upper height, a separated multiplicative filtration puts beta\_1 in degree at least two and vanishes beyond a dimension-controlled bound, killing that same power. Completion and naturality identify a nonzero lower-localized class in the unit's kernel, obstructing a retraction.

\subsection{Scope relative to this problem}
Manuscript 1 (family 318, 2026-09-27): The primary manuscript establishes nonzero sphere-kernel classes refuting the canonical weak chromatic splitting retraction at specified bad-prime heights. The bad-prime sphere kernels at n=p and n=p+1 give specific obstruction data for the overlap but not the requested systematic homotopy-type description of the entire bad-prime fiber.

\subsection{Formal verification}
Family 318 has no Lean scope page in this snapshot; no formalization of these manuscript results is listed.

This site has not executed the Lean files or checked their dependencies. These descriptions report the scope of the linked declarations, rather than a new verification of the proofs.

\subsection{Status caveat}
OpenAI attributes these results to an internal model. Its release notice cautions that results lacking full formalization can contain mistakes. This site has not independently assessed the proofs; a solved label records the scope of the reported claim.

\subsection{Sources}
\begin{itemize}
\item OpenAI, \emph{Counterexamples to weak chromatic splitting: sphere kernels and descent exponents}, 2026-09-27; citation key \texttt{OAI:Counterexamples-to-weak-chromatic-splitting-sphere-kernels-and-descent-exponents-September-27-2026}. \href{https://github.com/openai/math/blob/adc7f1241b42e322a6451854ab7e4b4c146bf78a/preprints/Counterexamples-to-weak-chromatic-splitting-sphere-kernels-and-descent-exponents-September-27-2026/paper.pdf}{PDF (GitHub)}; \href{https://github.com/openai/math/raw/adc7f1241b42e322a6451854ab7e4b4c146bf78a/preprints/Counterexamples-to-weak-chromatic-splitting-sphere-kernels-and-descent-exponents-September-27-2026/paper.pdf}{PDF (direct)}; \href{https://github.com/openai/math/blob/adc7f1241b42e322a6451854ab7e4b4c146bf78a/preprints/Counterexamples-to-weak-chromatic-splitting-sphere-kernels-and-descent-exponents-September-27-2026/README.md}{Manuscript page and citation}; \href{https://github.com/openai/math/blob/main/preprints/Counterexamples-to-weak-chromatic-splitting-sphere-kernels-and-descent-exponents-September-27-2026/paper.pdf}{Latest version}.
\item \href{https://github.com/openai/math/blob/adc7f1241b42e322a6451854ab7e4b4c146bf78a/CONTENTS.md}{OpenAI catalogue entry}.
\end{itemize}
\end{document}
